\chapter{Results}

All numbers below come from a single committed run on the i7-14700K, pinned to a
P core, at a software measured clock of about \SI{5.3}{\giga\hertz}.

\section{Cache latency}

The pointer chase recovers a clean four plateau staircase
(Figure~\ref{fig:latency}). L1 data latency is about \SI{0.94}{\nano\second},
which at the measured clock is 5.0 cycles, exactly the Raptor Cove L1 load use
latency. L2 is about 16 cycles, the L3 region climbs through the tens of cycles,
and DRAM settles near \SI{93}{\nano\second}. Table~\ref{tab:latency} lists the
detected plateaus.

\begin{figure}[h]
\centering
\includegraphics[width=\linewidth]{latency_plateau.png}
\caption{Load use latency against working set, with detected plateau means and
the sysfs cache capacities marked. The step to DRAM begins about one octave
below the nominal 33 MB L3, an effect of TLB reach and associativity under
random access.}
\label{fig:latency}
\end{figure}

\begin{table}[h]
\centering
\caption{Detected latency plateaus, cross checked against sysfs.}
\label{tab:latency}
\begin{tabular}{llrr}
\hline
Region & Working set & Latency (ns) & Cycles \\
\hline
L1D & 4 to 32 KiB & 0.94 & 5.0 \\
L2 & 64 to 2048 KiB & 4.29 & 22.9 \\
L3 & 4096 to 8192 KiB & 20.03 & 107.0 \\
DRAM & 16384 to 524288 KiB & 93.38 & 498.8 \\
\hline
\end{tabular}

\end{table}

The change point detector places the L1 boundary within \num{0.08} octaves of
the 48 KB L1 capacity and the L2 boundary within \num{0.50} octaves of the 2 MB
L2. The DRAM onset is flagged about \num{1.5} octaves below the nominal L3
capacity, which the discussion addresses.

\section{Bandwidth}

Single thread bandwidth by variant is in Table~\ref{tab:bandwidth} and
Figure~\ref{fig:variants}. The non temporal variant wins on Add and Triad, where
avoiding the read for ownership pays off, while the scalar and vectorized
variants are competitive on Copy. The thread scaling curve
(Figure~\ref{fig:scaling}) saturates near \SI{78}{\giga\byte\per\second} by four
to five threads, about 87 percent of the DDR5-5600 dual channel theoretical
\SI{89.6}{\giga\byte\per\second}. Beyond that point additional P cores, and then
the E cores, add no bandwidth: the reported knee is the bandwidth wall, not a P
to E transition.

\begin{table}[h]
\centering
\caption{Single thread STREAM bandwidth by variant, GB/s.}
\label{tab:bandwidth}
\begin{tabular}{lrrrr}
\hline
Variant & Copy & Scale & Add & Triad \\
\hline
Scalar & 43.9 & 23.1 & 27.1 & 28.7 \\
Auto vectorized & 46.1 & 26.3 & 31.8 & 31.4 \\
AVX2 NT stores & 36.7 & 36.8 & 41.3 & 40.8 \\
\hline
\end{tabular}

\end{table}

\begin{figure}[h]
\centering
\includegraphics[width=0.85\linewidth]{variant_comparison.png}
\caption{Single thread bandwidth by kernel and variant.}
\label{fig:variants}
\end{figure}

\begin{figure}[h]
\centering
\includegraphics[width=0.85\linewidth]{bandwidth_scaling.png}
\caption{Bandwidth scaling, P cores first then E cores. The curve saturates well
before the cores are exhausted.}
\label{fig:scaling}
\end{figure}

\section{Peak FLOPS}

Table~\ref{tab:flops} compares measured throughput to the ceiling computed from
the same measured clock. One P core reaches about \SI{170}{GFLOPS}, which is 99
percent of its ceiling; matching the ceiling to within one percent is the
expected result for a saturating ten chain loop and confirms both the clock and
the FLOP counting. All 28 logical CPUs together reach about \SI{1900}{GFLOPS}.

\begin{table}[h]
\centering
\caption{Peak single precision throughput against the computed ceiling.}
\label{tab:flops}
\begin{tabular}{lrrr}
\hline
Configuration & GFLOPS & Ceiling & Efficiency \\
\hline
Single P core & 169.7 & 172.4 & 98\% \\
All core (28 threads) & 1878 & (aggregate) & \\
\hline
\end{tabular}

\end{table}

\section{Roofline and tile prediction}

Figure~\ref{fig:roofline} places the suite's kernels on the measured roofline.
STREAM Triad sits on the memory bound slope. The single core FMA loop sits on the
single core ceiling. The naive blocked GEMM, though high in arithmetic intensity
and therefore compute bound in principle, sits well below the single core
ceiling: the gap is the cost of not packing.

Table~\ref{tab:tile} compares the BLIS analytical tile prediction to the
empirical best from the GEMM sweep. They disagree: the model predicts a large
$K_c$, the naive kernel prefers a much smaller one. This miss is a finding, not a
failure, and the discussion explains it.

\begin{figure}[h]
\centering
\includegraphics[width=\linewidth]{roofline.png}
\caption{Measured roofline with all core and single core compute ceilings.}
\label{fig:roofline}
\end{figure}

\begin{table}[h]
\centering
\caption{BLIS analytical prediction versus empirical best tile.}
\label{tab:tile}
\begin{tabular}{lrrr}
\hline
Source & Mc & Kc & Nc \\
\hline
BLIS analytical & 704 & 696 & 11296 \\
Empirical best & 384 & 64 & 1536 \\
\hline
\end{tabular}
% Kc prediction is a miss of 3.4 octaves in Kc.

\end{table}
